\paragraph{能量谱的热力学完成：大偏差、微正则熵与多分形读出}\label{par:pom-multiplicity-energy-ldp-multifractal}
沿用定义 \ref{def:pom-multiplicity-composition-partition} 的有序分解集合 \(\mathcal{C}_L\)、权重
\(
w(\boldsymbol{\ell})=\prod_i F_{\ell_i+2}
\)
与定义 \ref{def:pom-multiplicity-induced-composition-measure} 的纤维权测度
\(
P_L=\mathbb{P}_L
\)。
记对数多重度的能量密度为
\[
E_L(\boldsymbol{\ell}):=\frac{1}{L}\log w(\boldsymbol{\ell})
\qquad(\boldsymbol{\ell}\in\mathcal{C}_L).
\]
由命题 \ref{prop:pom-path-component-multiplicity-refinement-monotone-extrema} 的端点界 \(F_{L+2}\le w(\boldsymbol{\ell})\le 2^L\) 与 \(L^{-1}\log F_{L+2}\to\log\varphi\)，可得可达能量闭包为
\[
E_L(\boldsymbol{\ell})\in\bigl[L^{-1}\log F_{L+2},\,\log 2\bigr],
\qquad
\mathrm{cl}\,\mathrm{Range}(E_L)\subseteq[\log\varphi,\log 2].
\]

\begin{theorem}[纤维权测度下对数多重度的精确大偏差原理：压力函数 \texorpdfstring{$\log\lambda(q)$}{log lambda(q)} 的 Legendre 谱]\label{thm:pom-multiplicity-energy-ldp-under-PL}
\leanverified{paper\_pom\_multiplicity\_energy\_ldp\_under\_pl}
令 \(E_L\) 如上，且令 \(\lambda(q)\) 为定义 \ref{def:pom-multiplicity-real-q-pressure} 的实参数矩压力。
则随机变量序列 \((E_L)_{L\ge 1}\) 在 \((P_L)_{L\ge 1}\) 下满足大偏差原理；其规范化累积母函数在 \(t>-1\) 上存在并具有闭式
\[
\boxed{\;
\Psi(t):=\lim_{L\to\infty}\frac{1}{L}\log\mathbb{E}_{P_L}\!\left[e^{tLE_L}\right]
=\log\lambda(1+t)-\log\lambda(1)\; }.
\]
对应速率函数为
\[
\boxed{\;
I(e)=\sup_{t>-1}\bigl\{t e-\Psi(t)\bigr\}\in[0,\infty]\; } ,
\]
并且 \(I\) 在内部区间 \((\log\varphi,\log 2)\) 上严格凸且实解析，且在该区间外取 \(I(e)=+\infty\)。
\par
等价地，取参数 \(q>0\) 并令
\[
e(q):=\frac{d}{dq}\log\lambda(q)
\quad(\text{存在且由命题 \ref{prop:pom-multiplicity-lambdaq-derivative-gibbs} 给出显式表达}),
\]
则在值域内部有参数化闭式
\[
\boxed{\;
I\bigl(e(q)\bigr)=(q-1)e(q)-\log\frac{\lambda(q)}{\lambda(1)}\; } .
\]
\end{theorem}
\begin{proof}[证明要点]
对任意 \(t>-1\)（记 \(q=1+t>0\)），有
\[
\mathbb{E}_{P_L}\!\left[e^{tLE_L}\right]
=\sum_{\boldsymbol{\ell}\in\mathcal{C}_L}\frac{w(\boldsymbol{\ell})}{Z_L^{(1)}}\,w(\boldsymbol{\ell})^{t}
=\frac{Z_L^{(q)}}{Z_L^{(1)}},
\qquad
Z_L^{(q)}:=\sum_{\boldsymbol{\ell}\in\mathcal{C}_L}w(\boldsymbol{\ell})^{q}.
\]
由串接结构与母函数恒等式 \( \sum_{L\ge 0}Z_L^{(q)}t^L=1/(1-W_q(t))\)（见 \eqref{eq:pom-multiplicity-composition-moment-hierarchy-gf}，其推导对实 \(q>0\) 逐字成立），并结合 \(W_q(\rho(q))=1\)、\(\partial_tW_q(\rho(q))>0\)，可得
\(
\lim_{L\to\infty}L^{-1}\log Z_L^{(q)}=\log\lambda(q)
\)。
从而 \(\Psi(t)=\log\lambda(1+t)-\log\lambda(1)\)。
\par
由于定理 \ref{thm:pom-multiplicity-real-q-pressure} 给出 \(q\mapsto\log\lambda(q)\) 在 \((0,\infty)\) 上实解析且严格凸，故 \(\Psi\) 在 \(t>-1\) 上可微且在 \(t=0\) 的邻域解析；由 G\"artner--Ellis 定理得到 \(E_L\) 的大偏差原理与 Legendre--Fenchel 速率函数。
参数化形式由 \(e(q)=\Psi'(q-1)\) 与 \(I(e)=t e-\Psi(t)\) 在驻点处取值直接推出。证毕。
\end{proof}

\begin{corollary}[边界能量的指数稀有性常数：\texorpdfstring{$\log(\lambda(1)/2)$}{log(lambda(1)/2)} 与 \texorpdfstring{$\log(\lambda(1)/\varphi)$}{log(lambda(1)/phi)}]\label{cor:pom-multiplicity-energy-ldp-boundary}
\leanverified{paper\_pom\_multiplicity\_energy\_ldp\_boundary}
令 \(\varphi=\frac{1+\sqrt5}{2}\)，并注意由定理 \ref{thm:pom-multiplicity-real-q-pressure} 的整数一致性
\(
\lambda(1)=\lambda_1=\frac{3+\sqrt{17}}{2}
\)（见命题 \ref{prop:pom-multiplicity-composition-partition-rational}）。
则定理 \ref{thm:pom-multiplicity-energy-ldp-under-PL} 的速率函数在能量可达闭包端点处存在有限边界值，并满足
\[
\boxed{\;
\lim_{e\uparrow \log 2} I(e)=\log\frac{\lambda(1)}{2},\qquad
\lim_{e\downarrow \log\varphi} I(e)=\log\frac{\lambda(1)}{\varphi}\; } .
\]
因此，“全碎裂态”（\(\boldsymbol{\ell}=\boldsymbol{1}*L\)）与“单块态”（\(\boldsymbol{\ell}=(L)\)）在 \(P_L\) 下均呈线性指数稀有，其指数常数分别为 \(\log(\lambda(1)/2)\) 与 \(\log(\lambda(1)/\varphi)\)。
\end{corollary}
\begin{proof}
由 \(\boldsymbol{1}*L\) 的权重 \(w(\boldsymbol{1}*L)=2^L\) 与定义 \(P_L(\boldsymbol{\ell})=w(\boldsymbol{\ell})/Z_L^{(1)}\)，得
\(
P_L(\boldsymbol{1}*L)=2^L/Z_L^{(1)}
\)，故
\[
-\lim_{L\to\infty}\frac1L\log P_L(\boldsymbol{1}*L)=\log\lambda(1)-\log 2=\log\frac{\lambda(1)}{2}.
\]
同理对 \(\boldsymbol{\ell}=(L)\) 有 \(w(\boldsymbol{\ell})=F_{L+2}\)，且 \(L^{-1}\log F_{L+2}\to\log\varphi\)，从而
\[
-\lim_{L\to\infty}\frac1L\log P_L(\boldsymbol{\ell}=(L))=\log\lambda(1)-\log\varphi=\log\frac{\lambda(1)}{\varphi}.
\]
由于两端点各由单一原子实现，边界值亦可视为 \(I\) 的端点延拓；与定理 \ref{thm:pom-multiplicity-energy-ldp-under-PL} 的良性速率函数结构相容。证毕。
\end{proof}

\begin{theorem}[微正则计数熵的存在与全闭式对偶]\label{thm:pom-multiplicity-microcanonical-entropy}
\leanverified{paper\_pom\_multiplicity\_microcanonical\_entropy}
对 \(e\in\RR\)、\(\delta>0\) 定义能量壳计数
\[
N_L(e,\delta):=\#\Bigl\{\boldsymbol{\ell}\in\mathcal{C}_L:\ |E_L(\boldsymbol{\ell})-e|<\delta\Bigr\}.
\]
则极限
\[
s(e):=\lim_{\delta\downarrow 0}\ \lim_{L\to\infty}\frac{1}{L}\log N_L(e,\delta)
\]
在 \(e\in(\log\varphi,\log 2)\) 上存在，并满足热力学对偶恒等式
\[
\boxed{\;
\log\lambda(q)=\sup_{e\in[\log\varphi,\log 2]}\bigl(s(e)+q e\bigr)\qquad(q\ge 0)\; } ,
\]
从而
\[
\boxed{\;
s(e)=\inf_{q\ge 0}\bigl(\log\lambda(q)-q e\bigr)\qquad(e\in[\log\varphi,\log 2])\; } .
\]
并且 \(s\) 在内部区间严格凹且实解析，且边界满足
\[
\boxed{\;
\lim_{e\downarrow \log\varphi}s(e)=\lim_{e\uparrow \log 2}s(e)=0\; } .
\]
\end{theorem}
\begin{proof}[证明要点]
由定义 \(Z_L^{(q)}=\sum_{\boldsymbol{\ell}\in\mathcal{C}_L}e^{qLE_L(\boldsymbol{\ell})}\) 与极限
\(
\lim_{L\to\infty}L^{-1}\log Z_L^{(q)}=\log\lambda(q)
\)
（同定理 \ref{thm:pom-multiplicity-energy-ldp-under-PL} 证明要点），对 \(E_L\) 的计数测度应用 Varadhan 型极大项原则可得
\(
\log\lambda(q)=\sup_e(s(e)+q e)
\)。
对 \(q\) 作 Legendre--Fenchel 对偶即得到 \(s(e)=\inf_{q\ge 0}(\log\lambda(q)-q e)\)。
严格凹性与解析性来自 \(q\mapsto\log\lambda(q)\) 的严格凸解析性（定理 \ref{thm:pom-multiplicity-real-q-pressure}）的对偶传递。
端点 \(s\to 0\) 由端点能量各仅由单一分解实现（命题 \ref{prop:pom-path-component-multiplicity-refinement-monotone-extrema}）给出。证毕。
\end{proof}

\begin{corollary}[Gibbs 身份：纤维权测度的能量大偏差与微正则熵的精确关系]\label{cor:pom-multiplicity-gibbs-identity}
\leanverified{paper\_pom\_multiplicity\_gibbs\_identity}
令 \(I\) 为定理 \ref{thm:pom-multiplicity-energy-ldp-under-PL} 的速率函数，\(s\) 为定理 \ref{thm:pom-multiplicity-microcanonical-entropy} 的微正则计数熵。
则在内部区间 \((\log\varphi,\log 2)\) 上成立严格恒等式
\[
\boxed{\;
I(e)=\log\lambda(1)-s(e)-e\; } ,
\]
并可连续延拓到端点（对应推论 \ref{cor:pom-multiplicity-energy-ldp-boundary} 的边界值）。
特别地，典型能量
\(
e_\star=\left.\frac{d}{dq}\log\lambda(q)\right|_{q=1}
\)
唯一满足支撑线条件
\(
s'(e_\star)=-1
\)，且
\(
s(e_\star)=\log\lambda(1)-e_\star
\)。
\end{corollary}
\begin{proof}
由 \(P_L(\boldsymbol{\ell})=e^{LE_L(\boldsymbol{\ell})}/Z_L^{(1)}\)，对能量壳概率有近似分解
\[
P_L\bigl(|E_L-e|<\delta\bigr)
=\sum_{|E_L-e|<\delta}\frac{e^{LE_L}}{Z_L^{(1)}}
=\frac{e^{Le+o(L)}}{Z_L^{(1)}}\,N_L(e,\delta).
\]
取对数、除以 \(L\)，再用 \(L^{-1}\log Z_L^{(1)}\to\log\lambda(1)\) 与 \(L^{-1}\log N_L(e,\delta)\to s(e)\) 即得
\(
I(e)=\log\lambda(1)-s(e)-e
\)。
支撑线条件由 \(s(e)=\inf_q(\log\lambda(q)-q e)\) 的驻点方程给出。证毕。
\end{proof}

\begin{proposition}[原子权重的多分形谱：\texorpdfstring{$\alpha$}{alpha} 区间与谱函数闭式]\label{prop:pom-multiplicity-atom-multifractal}
\leanverified{paper\_pom\_multiplicity\_atom\_multifractal}
令单原子负对数概率密度为
\[
\alpha_L(\boldsymbol{\ell}):=-\frac{1}{L}\log P_L(\boldsymbol{\ell})\qquad(\boldsymbol{\ell}\in\mathcal{C}_L).
\]
则有一致渐近
\[
\alpha_L(\boldsymbol{\ell})=\log\lambda(1)-E_L(\boldsymbol{\ell})+o(1),
\]
从而其可达极限区间为
\[
\boxed{\;
\alpha\in\Bigl[\log\frac{\lambda(1)}{2},\ \log\frac{\lambda(1)}{\varphi}\Bigr]\; } .
\]
并且对任意内部 \(\alpha\)（令 \(e=\log\lambda(1)-\alpha\)），原子指数谱满足
\[
\boxed{\;
\lim_{\delta\downarrow 0}\ \lim_{L\to\infty}\frac{1}{L}\log
\#\Bigl\{\boldsymbol{\ell}\in\mathcal{C}_L:\ |\alpha_L(\boldsymbol{\ell})-\alpha|<\delta\Bigr\}
=f(\alpha):=s\bigl(\log\lambda(1)-\alpha\bigr)\; } ,
\]
其中 \(s\) 为定理 \ref{thm:pom-multiplicity-microcanonical-entropy} 的微正则计数熵。
因此 \(f\) 在内部严格凹且实解析，并满足端点值 \(f(\log(\lambda(1)/2))=f(\log(\lambda(1)/\varphi))=0\)。
\end{proposition}
\begin{proof}[证明要点]
由 \(P_L(\boldsymbol{\ell})=w(\boldsymbol{\ell})/Z_L^{(1)}\) 与 \(L^{-1}\log Z_L^{(1)}\to\log\lambda(1)\)，得 \(\alpha_L=\log\lambda(1)-E_L+o(1)\)。
计数谱结论由 \(\alpha\leftrightarrow e\) 的双射与定理 \ref{thm:pom-multiplicity-microcanonical-entropy} 的 \(s(e)\) 直接推出。证毕。
\end{proof}

\begin{proposition}[连续 Rényi 熵率/维谱：整数结果的全实数延拓]\label{prop:pom-multiplicity-renyi-continuous}
\leanverified{paper\_pom\_multiplicity\_renyi\_continuous}
对实数 \(q>0,q\ne 1\) 定义 Rényi 熵
\[
H_q(P_L):=\frac{1}{1-q}\log_2\sum_{\boldsymbol{\ell}\in\mathcal{C}_L}P_L(\boldsymbol{\ell})^{q}.
\]
则极限熵率存在并满足闭式
\[
\boxed{\;
\lim_{L\to\infty}\frac{1}{L}H_q(P_L)
=\frac{1}{1-q}\Bigl(\log_2\lambda(q)-q\log_2\lambda(1)\Bigr)\; } .
\]
并且 Shannon 熵率可由可导极限给出：
\[
\boxed{\;
\lim_{L\to\infty}\frac{1}{L}H(P_L)
=\log_2\lambda(1)-\left.\frac{d}{dq}\log_2\lambda(q)\right|_{q=1}\; } .
\]
\end{proposition}
\begin{proof}
由 \(P_L(\boldsymbol{\ell})=w(\boldsymbol{\ell})/Z_L^{(1)}\) 得
\(
\sum P_L^q=Z_L^{(q)}/(Z_L^{(1)})^q
\)。
对数除以 \(L\) 并用 \(L^{-1}\log Z_L^{(q)}\to\log\lambda(q)\) 得到 Rényi 熵率闭式；令 \(q\to 1\) 并用 \(\log\lambda\) 的可微性给出 Shannon 熵率公式。证毕。
\end{proof}

\begin{theorem}[对数多重度的中心极限定律与方差常数：\texorpdfstring{$\log\lambda$}{log lambda} 的二阶导数]\label{thm:pom-multiplicity-logw-clt}
在 \(P_L\) 下把 \(\log w(\boldsymbol{\ell})\) 视为随机变量。
则存在常数
\[
e_\star:=\left.\frac{d}{dq}\log\lambda(q)\right|_{q=1},
\qquad
\sigma_\star^2:=\left.\frac{d^2}{dq^2}\log\lambda(q)\right|_{q=1}\in(0,\infty),
\]
使得中心极限定律成立：
\[
\boxed{\;
\frac{\log w(\boldsymbol{\ell})-e_\star L}{\sqrt{L}}\ \Longrightarrow\ \mathcal{N}(0,\sigma_\star^2)\; } .
\]
等价地，
\(
\sqrt{L}\,(E_L-e_\star)\Longrightarrow \mathcal{N}(0,\sigma_\star^2)
\)。
\end{theorem}
\begin{proof}[证明要点]
由定理 \ref{thm:pom-multiplicity-energy-ldp-under-PL}，\(\Psi(t)=\log\lambda(1+t)-\log\lambda(1)\) 在 \(t=0\) 的邻域解析且严格凸；
因此可应用标准的 quasi-power/解析扰动型中心极限定理（或由 Cramér 型方法对 \(\Psi\) 作二阶展开并 Fourier 反演），得到 \(\log w\) 的 CLT。
方差正性由严格凸性推出。证毕。
\end{proof}

\begin{proposition}[中偏差原理：二次型速率由 \texorpdfstring{$\sigma_\star^2$}{sigma^2} 完全决定]\label{prop:pom-multiplicity-logw-mdp}
取任意序列 \(a_L\to\infty\) 且 \(a_L=o(\sqrt{L})\)。
则在 \(P_L\) 下成立中偏差原理：对任意 \(x\in\RR\)，
\[
\boxed{\;
\lim_{L\to\infty}\frac{1}{a_L^2}\log
\mathbb{P}\!\left(\frac{\log w(\boldsymbol{\ell})-e_\star L}{a_L\sqrt{L}}\ge x\right)
=-\frac{x^2}{2\sigma_\star^2}\; } ,
\]
且相同二次型速率对左尾亦成立。
\end{proposition}
\begin{proof}[证明要点]
由 \(\Psi\) 在零点邻域的解析性与二阶展开 \(\Psi(t)=e_\star t+\frac12\sigma_\star^2 t^2+O(t^3)\)，对缩放 \(t=\theta\,a_L/\sqrt{L}\) 应用 G\"artner--Ellis 的中偏差版本可得结论。证毕。
\end{proof}

\begin{theorem}[均匀分解测度下的能量谱与速率：\texorpdfstring{$\log\lambda(q)-\log 2$}{log lambda(q)-log 2} 的 Legendre 对偶]\label{thm:pom-multiplicity-energy-ldp-uniform}
令 \(U_L\) 为 \(\mathcal{C}_L\) 上的均匀测度（故 \(U_L(\boldsymbol{\ell})=2^{-(L-1)}\)）。
则 \((E_L)_{L\ge 1}\) 在 \((U_L)\) 下满足大偏差原理，其累积母函数为
\[
\boxed{\;
\Phi(q):=\lim_{L\to\infty}\frac{1}{L}\log\mathbb{E}_{U_L}\!\left[e^{qLE_L}\right]
=\log\lambda(q)-\log 2\qquad(q\ge 0)\; } ,
\]
速率函数为
\[
\boxed{\;
J(e)=\sup_{q\ge 0}\bigl\{q e-\Phi(q)\bigr\}\; } .
\]
并且端点满足
\[
\boxed{\;
\lim_{e\uparrow\log 2}J(e)=\lim_{e\downarrow\log\varphi}J(e)=\log 2\; } ,
\]
对应两端各仅有单原子实现的计数稀有性（\(|\mathcal{C}_L|=2^{L-1}\)）。
\end{theorem}
\begin{proof}[证明要点]
由 \(\mathbb{E}_{U_L}[e^{qLE_L}]=2^{-(L-1)}Z_L^{(q)}\) 与 \(L^{-1}\log Z_L^{(q)}\to\log\lambda(q)\) 得 \(\Phi(q)=\log\lambda(q)-\log 2\)。
再由 \(\Phi\) 的可微性与严格凸性应用 G\"artner--Ellis 定理得到 LDP 与 \(J\) 的 Legendre 形式。
端点值由端点事件为单原子且其均匀概率为 \(2^{-(L-1)}\) 直接得到。证毕。
\end{proof}

\begin{theorem}[正则—微正则等价：条件能量壳的局部极限为 \texorpdfstring{$q$}{q}-倾斜 i.i.d. 过程]\label{thm:pom-multiplicity-ensemble-equivalence}
对任意 \(q>0\) 定义 \(q\)-倾斜纤维权测度 \(P_L^{(q)}\)（见命题 \ref{prop:pom-multiplicity-low-temperature-gibbs}）：
\[
P_L^{(q)}(\boldsymbol{\ell})=\frac{w(\boldsymbol{\ell})^{q}}{Z_L^{(q)}},
\qquad
Z_L^{(q)}=\sum_{\boldsymbol{\ell}\in\mathcal{C}_L}w(\boldsymbol{\ell})^{q}.
\]
并令 \(p_q(k)=F_{k+2}^{\,q}\rho(q)^k\)（\(k\ge 1\)）为 \(\rho(q)=\lambda(q)^{-1}\) 所诱导的 i.i.d.\ 部件长度分布（见 \eqref{eq:pom-multiplicity-low-temp-iid-law}）。
记 \(e(q)=\frac{d}{dq}\log\lambda(q)\)。
\par
则对任意固定窗口 \(m\ge 1\) 与任意有界连续函数 \(\mathcal{F}\)（依赖前 \(m\) 个部件），成立 Gibbs 条件原理：
\[
\boxed{\;
\lim_{\delta\downarrow 0}\ \lim_{L\to\infty}
\mathbb{E}_{P_L}\!\left[\mathcal{F}(\ell_1,\dots,\ell_m)\ \middle|\ |E_L-e(q)|<\delta\right]
=\mathbb{E}_{p_q^{\otimes m}}[\mathcal{F}]\; } .
\]
因此，能量密度的微正则限制在局部上等价于 \(q\)-倾斜的正则系综；其中 \(q=1\) 的特例给出纤维权测度 \(P_L\) 的局部 i.i.d.\ 极限律。
\end{theorem}
\begin{proof}[证明要点]
由 \(P_L=P_L^{(1)}\) 且
\(
\frac{dP_L^{(q)}}{dP_L}(\boldsymbol{\ell})\propto e^{(q-1)LE_L(\boldsymbol{\ell})}
\)，
\(P_L^{(q)}\) 是 \(P_L\) 的指数倾斜。
对 \(P_L\) 下的稀有事件 \(|E_L-e(q)|<\delta\) 应用标准的 Gibbs 条件原理（由定理 \ref{thm:pom-multiplicity-energy-ldp-under-PL} 的严格凸速率函数保证唯一倾斜参数），可得该条件下有限窗口分布与 \(P_L^{(q)}\) 的有限窗口分布渐近相同。
再由命题 \ref{prop:pom-multiplicity-low-temperature-gibbs}，在 \(P_L^{(q)}\) 下有限窗口极限为 \(p_q^{\otimes m}\)，合成即得结论。证毕。
\end{proof}

\endinput
